\documentclass[10pt]{article}
% Math Packages
\usepackage{amsmath, mathtools}
\usepackage{amssymb}
\usepackage{amsthm}
\usepackage{amsfonts}
\usepackage{bbm}
\usepackage{breqn}
\usepackage[margin=1in]{geometry}
\usepackage{graphicx}
\usepackage{tikz}
\usetikzlibrary{arrows.meta}
\usetikzlibrary{calc}
\usepackage{forest}
\usepackage{tikz-qtree}
\graphicspath{ {./images/} }
\usepackage{hyperref}
\usepackage[capitalize]{cleveref}
\usepackage[shortlabels]{enumitem}
\usetikzlibrary{arrows,matrix,positioning}
\usepackage{multicol}

\usepackage{listings}
\usepackage{xcolor}

\crefname{problem}{Problem}{Problems}
\Crefname{problem}{Problem}{Problems}

\lstset{
  language=R,
  basicstyle=\ttfamily\footnotesize,
  keywordstyle=\color{blue},
  commentstyle=\color{gray},
  stringstyle=\color{red!70!black},
  framesep=2pt,
  framexleftmargin=0pt, 
  frame=single,
  breaklines=false,
  showstringspaces=false
}
% for the pipe symbol
\usepackage[T1]{fontenc}

% Colorful Notes
\usepackage{color} \definecolor{Red}{rgb}{1,0,0} \definecolor{Blue}{rgb}{0,0,1}
\definecolor{Purple}{rgb}{.5,0,.5} \def\red{\color{Red}} \def\blue{\color{Blue}}
\def\gray{\color{gray}} \def\purple{\color{Purple}}
\newcommand{\rnote}[1]{{\red [#1]}} % \rnote{foo} gives '[foo]' in red
\newcommand{\pnote}[1]{{\purple [#1]}} % \pnote{foo} gives '[foo]' in purple
\newcommand{\bnote}[1]{{\blue #1}} % \bnote{foo} gives 'foo' in blue
\newcommand{\gnote}[1]{{\gray #1}} % \gnote{foo} gives 'foo' in gray
\newcommand{\Max}[1]{{\purple [#1]}} % \bnote{foo} then 'foo' is blue

% Claim numbering (the counter restarts after each proof environment)
\newcounter{claimcount}
\setcounter{claimcount}{0}
\newenvironment{claim}{\refstepcounter{claimcount}\par\addvspace{\medskipamount}\noindent\textbf{Claim \arabic{claimcount}:}}{}
\usepackage{etoolbox}
\AtBeginEnvironment{proof}{\setcounter{claimcount}{0}}
\newenvironment{claimproof}{\par\addvspace{\medskipamount}\noindent\textit{Proof of Claim  \arabic{claimcount}.}}{\hfill\ensuremath{\qedsymbol} \tiny{Claim}

  \medskip}
% Add claim support to cleverref
\crefname{claimcount}{Claim}{Claims}

% Math Environments
\newtheorem{theorem}{Theorem}
\newtheorem{assumption}[theorem]{Assumption}
\newtheorem{lemma}[theorem]{Lemma}
\newtheorem{proposition}[theorem]{Proposition}
\newtheorem{corollary}[theorem]{Corollary}
\newtheorem{question}[theorem]{Question}
\theoremstyle{definition}
\newtheorem*{definition}{Definition}
\newtheorem{remark}[theorem]{Remark}
\newtheorem{example}[theorem]{Example}
\newtheorem{notation}[theorem]{Notation}
\newtheorem{problem}[theorem]{Problem}

% Matrices and Column Vectors. 
\usepackage{stackengine}
\setstackgap{L}{1.0\normalbaselineskip}
\usepackage{tabstackengine}
\setstackEOL{;}% row separator
\setstackTAB{,}% column separator
\setstacktabbedgap{1ex}% inter-column gap
\setstackgap{L}{1.5\normalbaselineskip}% inter-row baselineskip
\let\nmatrix\bracketMatrixstack  %Usage: \nmatrix{1,2,3\4,5,6}
\newcommand\cv[1]{\setstackEOL{,}\bracketMatrixstack{#1}} %usage: \cv{1,2,3}

% Custom Math Coqmmands
\newcommand{\vt}{\vskip 5mm} % vertical space
\newcommand{\fl}{\noindent\textbf} % first line
\newcommand{\Fl}{\vt\noindent\textbf} % first line with space above
\newcommand{\norm}[1]{\left\lVert#1\right\rVert} % norm
\newcommand{\pnorm}[1]{\left\lVert#1\right\rVert_p} % p-norm
\newcommand{\qnorm}[1]{\left\lVert#1\right\rVert_q} % q-norm
\newcommand{\1}[1]{\textbf{1}_{\left[#1\right]}} % indicator function 
\def\limn{\lim_{n\to\infty}} % shortcut for lim as n-> infinity
\def\sumn{\sum_{n=1}^{\infty}} % shortcut for sum from n=1 to infinity
\def\sumkn{\sum_{k=1}^{n}} % shortcut for sum from k=1 to n
\def\sumin{\sum_{i=1}^{n}} % shortcut for sum from i=1 to n
\def\SAs{\sigma\text{-algebras}} % shortcut for $\sigma$-algebras
\def\SA{\sigma\text{-algebra}} % shortcut for $\sigma$-algebra
\def\Ft{\mathcal{F}_t} % time-indexed sigma-algebra (t)
\def\Fs{\mathcal{F}_s} % time-indexed sigma-algebra (s)
\def\F{\mathcal{F}} % sigma-algebra
\def\G{\mathcal{G}} % sigma-algebra
\def\R{\mathbb{R}} % Real numbers
\def\N{\mathbb{N}} % Natural numbers
\def\Z{\mathbb{Z}} % Integers
\def\E{\mathbb{E}} % Expectation
\def\P{\mathbb{P}} % Probability
\def\Q{\mathbb{Q}} % Q probability
\def\dist{\text{dist}} %Text 'dist' for things like 'dist(x,y)'
\newcommand{\indep}{\perp \!\!\! \perp}  %independence symbol
\def\Var{\mathrm{Var}} % Variance
\def\tr{\mathrm{tr}} % trace

% Brackets and Parentheses
\def\[{\left [}
    \def\]{\right ]}
% \def\({\left (}
%   \def\){\right )}

\usepackage{color}
\definecolor{Red}{rgb}{1,0,0}
\definecolor{Blue}{rgb}{0,0,1}
\definecolor{Purple}{rgb}{.5,0,.5}
\def\red{\color{Red}}
\def\blue{\color{Blue}}
\def\gray{\color{gray}}
\def\purple{\color{Purple}}
\definecolor{RoyalBlue}{cmyk}{1, 0.50, 0, 0}
\newcommand{\dempfcolor}[1]{{\color{RoyalBlue}#1}} 
\newcommand{\demph}[1]{\textcolor{RoyalBlue}{\textbf{\slshape #1}}} % Slanted

\usepackage{emoji}
% comment exactly one of the following line to show / hide the solutions
% \newcommand{\solution}[1]{{\purple #1}} % uncomment to show the solutions
\newcommand{\solution}[1]{} % uncomment to hide the solutions




\begin{document}
\begin{center}
  \section*{Math 472: Homework 08}
  \textit{Due Monday, April 20}
\end{center}


\begin{problem}[Exercise 10.6]
  \label{problem:exercise-10.6}
  We are interested in testing whether a coin is balanced based on the number
  of heads $Y$ in 36 flips; that is, $H_{0}:p=.5$ vs $H_{a}:p\neq 5$. If we
  use the rejection region $|y-18| \geq 4$, what is
  \begin{enumerate}[(a)]
    \item the value of $\alpha$?
    \item the value of $\beta$ if $p=.7$?
  \end{enumerate}
\end{problem}


\begin{problem}
  Please answer ``true'' or ``false'' to the following questions.
  \begin{enumerate}[(a)]
    \item The level of the test computed in \cref{problem:exercise-10.6}(a) is
    the probability that $H_{0}$ is true.
    \item The value of $\beta$ in \cref{problem:exercise-10.6}(b) is the
    probability that $H_{a}$ is true.
    \item In \cref{problem:exercise-10.6}(b), $\beta$ was computed assuming
    the null hypothesis was false.
    \item If $\beta$ was computed when $p=.55$, the value would be larger than
    the value of $\beta$ obtained in \cref{problem:exercise-10.6}(b).
    \item The probability that the test mistakenly rejects $H_{0}$ is $\beta$. 
    \item Suppose the rejection region was changed to $\left| y-18 \right|
    \geq 2$.
    \begin{enumerate}[(i)]
      \item This RR would lead to rejecting the null hypothesis more often
      than the RR used in \cref{problem:exercise-10.6}.
      \item If $\alpha$ was comptued using this new RR, the value would be
      larger than the value obtained in \cref{problem:exercise-10.6}(b).
      \item If $\beta$ was computed when $p=.7$ and using this new RR, the
      value would be alrger than the valued obtained in
      \cref{problem:exercise-10.6}(b).
    \end{enumerate}
  \end{enumerate}
\end{problem}


\begin{problem}[Exercise 10.106]
  % I love this problem.
  A survey of voter sentiment was conducted in four political wards to compare
  the fraction of voters favoring candidate Jones. Random samples of 200
  voters were polled in each of the four wards, with the results as shown in
  the following table:
  \begin{center}
    \begin{tabular}{lccccc}
      & \multicolumn{4}{c}{Ward} & \\
      \cline{2-5}
      Opinion & 1 & 2 & 3 & 4 & Total \\
      \hline
      Favor A & 76 & 53 & 59 & 48 & 236 \\
      Do not favor A & 124 & 147 & 141 & 152 & 564 \\
      \hline
      Total & 200 & 200 & 200 & 200 & 800 \\
    \end{tabular}
  \end{center}
  The number of voters favoring Jones in the four samples can be regarded as
  four independent binomal random variables $X_{1},X_{2},X_{3},$ and $X_{4}$,
  of the form $X_{i}\sim \text{Binomial}(200,p_{i})$ for $i=1,2,3,4$, where
  $p_{i}$ is the proportion of voters in ward $i$ who support Jones.
  
  In the following parts we will construct a likelihood ratio test of the
  hypothesis that the fraction of voters favoring candidate Jones are the same
  in all four wards.
  \begin{enumerate}[(a)]
    \item State the null and alternative hypotheses.
    \item Describe the null hypothesis geometrically. What is the dimension of
    the null hypothesis?
    \item Write down the joint likelihood of $X_{1},X_{2},X_{3}$ and $X_{4}$
    and find its maximum.
    \item Write down the joint likelihood of $X_{1},X_{2},X_{3}$ and $X_{4}$
    when $H_{0}$ is true, and find its maximum.
    \item Using your answer to the previous two parts, compute $\Lambda$ as it
    was defined in lecture:
    \begin{equation*}
      \Lambda 
      = -2\log\left( \frac{\sup_{\theta\in \Theta_{0}}L(\theta)}
        {\sup_{\theta\in \Theta}L(\theta)}  \right) 
    \end{equation*}
    \item Compute a $p$-value for the data, and interpret your result in
    words.
  \end{enumerate}
\end{problem}


\begin{problem}[Exercise 10.88]
  Refer to problem 4 in Homework 8. Find the power of the test for each
  alternative in $(a)-(d)$:
  \begin{enumerate}[(a)]
    \item $H_{a}:p=.4 $
    \item $H_{a}:p=.5 $
    \item $H_{a}:p=.6 $
    \item $H_{a}:p=.7 $
    \item Sketch a graph of the power function.
  \end{enumerate}
\end{problem}


\begin{problem}[Method of Moments - exercise 9.69]
  Let $X$ be the proportion of allotted time that a randomly-selected student
  spends working on an exam. Suppose that $X$ is a random variable with pdf
  \begin{equation*}
    f(x) = \left\{ \begin{array}{l@{\quad:\quad}l} (\theta+1)x^{\theta}
        &0 \leq x \leq 1 \\ 0& \text{otherwise} \end{array}\right.
  \end{equation*}
  where $\theta>-1$ is an unknown parameter. A random sample of 10 students
  yields data
  \begin{gather*}
    x_{1} = .92,\ x_{2} = .79,\ x_{3} = .90,\ x_{4}=.65,\ x_{5}=.86\ \\
    x_{6}=.47,\ x_{7}=.73,\ x_{8}=.97,\ x_{9}=.94,\ x_{10}=.77
  \end{gather*}
  \begin{enumerate}[(a)]
    \item Use the method of moments to obtain an estimator of $\theta$ and
    then compute the estimate for this data.
    \item Use the factorization criterion to show that the statistic
    $T(X_{1},\ldots,X_{n})= - \sum_{i=1}^{n} \log(Y_{i})$ is sufficient for
    estimating $\theta$. 
    \item Is the method of moments estimator a function of $T$? What
    implications does this have?
  \end{enumerate}
\end{problem}


\begin{problem}[Exercise 9.19]
  Let $Y_{1},\ldots,Y_{n}$ be a random sample from the pdf
  \begin{equation*}
    f(x) = \theta y^{\theta-1}\mathbf{1}_{\left[ 0<y<1 \right]},
  \end{equation*}
  where $\theta>0$. Show that $\overline{Y}$ is a consistent estimator of
  $\frac{\theta}{\theta+1}$.
\end{problem}


\end{document}
% let the jangle of the scramsax be the only sound I hear